\section{Обзор существующих решений}
Перед выбором метода рекомендаций необходимо проанализировать существующие программы, выполняющие аналогичную задачу, и сравнить их.
Сравнение выполнено по следующим критериям:
\begin{itemize}
    \item наличие возможности запрашивать рекомендации рецептов -- K1;
    \item количество различных оценок, присваиваемых рецепту -- К2;
    \item наличие возможности указать все доступные пользователю ингредиенты -- K3;
    \item наличие возможности самостоятельного поиска рецепта по базе данных -- К4;
    \item наличие возможности получения рекомендаций без интернета -- К5;
    \item наличие возможности синхронизации данных о пользователе на разных устройствах -- K6;
    \item наличие возможности использования полностью на русском языке -- К7.
\end{itemize}

Рассмотрены популярные приложения, предназначенные для подбора рецептов блюд.
Приложение считается популярным, если количество его пользователей больше миллиона.
В результате выбраны следующие приложения:
\begin{itemize}
    \item Tasty~\cite{app_tasty};
    \item Простые рецепты на каждый день~\cite{app_easyrecipes};
    \item Food.ru пошаговые рецепты~\cite{app_foodru};
    \item Поваренная книга рецептов~\cite{app_cookbook};
    \item Простые рецепты~\cite{app_simple_recipes}.
\end{itemize}

\begin{table}[h]
\centering
\caption{ \centering Сравнение существующих приложений по критериям К1--К7}
\label{tab:overview}
\begin{tabular}{|l|c|c|c|c|c|c|c|c}
\hline
\textbf{Приложение} & \textbf{К1} & \textbf{К2} & \textbf{К3} & \textbf{К4} & \textbf{К5} & \textbf{К6} & \textbf{К7} \\
\hline
Tasty & + & 1 & -- & + & -- & + & + \\
\hline
Простые рецепты на каждый день & + & 5 & -- & + & + & + & + \\
\hline
Food.ru пошаговые рецепты & + & 5 & -- & + & -- & + & + \\
\hline
Поваренная книга рецептов & -- & 0 & -- & + & -- & -- & -- \\
\hline
Простые рецепты & -- & 0 & -- & + & -- & -- & + \\
\hline
\end{tabular}
\end{table}

Из результатов сравнения следует, что разрабатываемое приложение должно удовлетворять всем сформулированным критериям.
Именно тогда оно сможет быть лучше каждого рассмотренного аналога относительно сформулированных критериев.